
Let $\Fl(n)$ denote the variety of complete flags in $\C^n$, and let
$\Fl(\bfr,n)=\Fl(r_1, \dots, r_k; n)$ be the variety of partial flags. These are homogeneous under the group
$\SL_n(\C)$, and the restriction of this action to the maximal torus
$T \subset \SL_n(\C)$ has finitely many fixed points, indexed
by a~quotient of the symmetric group $S_n$.
Denote
by $\QK_T(\Fl(\bfr,n))$ the (equivariant, small) quantum K
ring associated to these varieties. This is an algebra over
$\K_T(\pt)[\![Q_1, \dots, Q_k]\!]$, and it has
a~$\K_T(\pt)[\![Q_1, \dots, Q_k]\!]$-basis given by Schubert classes
$\cO^w$ indexed by the torus fixed points. The quantum K multiplication
was defined by Givental and Lee~\cite{givental:onwdvv,lee:QK} in
terms of $3$-point, genus~$0$, $\K$-theoretic Gromov--Witten (KGW) invariants.
Denote by
\[ 0=\cS_0 \subset \cS_1 \subset \dots \subset \cS_k \subset \cS_{k+1}=\C^n \]
the sequence of tautological bundles in $\Fl(r_1, \dots, r_k; n)$; thus
$\mathrm{rank}(\cS_i)=r_i$ for $0\leq i\leq k+1$ with $r_0=0$ and $r_{k+1}=n$.

While the computational foundations of the quantum K rings
of (cominuscule) Grassmannians have been studied for some
time now (see, e.g.,~\cite{BCMP:qkpos,BCMP:qkchev,buch.m:qk,chaput.perrin:rationality,Gorbounov:2014,sinha2024quantum}),
